\chapter{จุลินทรีย์ในการหมักและการจัดการของเสียทางการเกษตร}
\label{chap:ch06}

% ==========================================================
% แผนบริหารการสอนประจำบทที่ 6
% ==========================================================
\section*{แผนบริหารการสอนประจำบทที่ 6}
\addcontentsline{toc}{section}{แผนบริหารการสอนประจำบทที่ 6}

\textbf{หัวข้อเนื้อหาประจำบท}
\begin{enumerate}
    \item ปริมาณและลักษณะของวัสดุเหลือใช้ทางการเกษตรในประเทศไทย
    \item ชีววิทยาและลำดับขั้นการสืบแทนของจุลินทรีย์ในการทำปุ๋ยหมัก (Composting)
    \item สภาวะที่เหมาะสมต่อการหมัก (อัตราส่วน C:N ความชื้น การถ่ายเทอากาศ และอุณหภูมิ)
    \item กระบวนการย่อยสลายแบบไม่ใช้ออกซิเจนและการผลิตก๊าซชีวภาพ (Biogas)
    \item นวัตกรรมการผลิตน้ำหมักชีวภาพและจุลินทรีย์สังเคราะห์แสง (PSB)
\end{enumerate}

\textbf{วัตถุประสงค์เชิงพฤติกรรม}
\begin{itemize}
    \item อธิบายกลไกการเปลี่ยนแปลงอุณหภูมิและการสืบแทนของจุลินทรีย์ในกองปุ๋ยหมักได้
    \item คำนวณอัตราส่วน C:N ของวัตถุดิบผสมสำหรับการทำปุ๋ยหมักได้อย่างถูกต้อง
    \item อธิบาย 4 ขั้นตอนชีวเคมีในการหมักก๊าซชีวภาพจากมูลสัตว์และของเสียเกษตรได้
    \item วางแผนและควบคุมกระบวนการหมักปุ๋ยอินทรีย์ให้ได้มาตรฐานความปลอดภัยได้
\end{itemize}

\textbf{กิจกรรมการเรียนการสอน}
\begin{itemize}
    \item การบรรยายและวิเคราะห์กราฟอุณหภูมิการหมักปุ๋ยในระดับอุตสาหกรรม
    \item ปฏิบัติการตั้งกองปุ๋ยหมักและตรวจวัดอุณหภูมิ ก๊าซคาร์บอนไดออกไซด์ และค่าความเป็นกรด-ด่าง
    \item การเพาะเลี้ยงแบคทีเรียสังเคราะห์แสงสีแดง (\textit{Rhodopseudomonas palustris})
\end{itemize}

\textbf{สื่อการเรียนการสอน}
\begin{itemize}
    \item เอกสารคำสอนบทที่ 6 และแผนภาพเวกเตอร์เส้นโค้งอุณหภูมิการหมัก
    \item เทอร์โมมิเตอร์ก้านยาวสำหรับวัดอุณหภูมิกองปุ๋ยหมัก
\end{itemize}

\textbf{การประเมินผล}
\begin{itemize}
    \item การประเมินผลการตรวจวัดคุณภาพปุ๋ยหมัก (ระดับการย่อยสลาย กลิ่น และสี)
    \item การทดสอบการคำนวณสูตรผสมวัตถุดิบปุ๋ยหมัก
\end{itemize}

\newpage

% ==========================================================
% เนื้อหาบทเรียน
% ==========================================================

\section{ชีววิทยาของการทำปุ๋ยหมักแบบใช้ออกซิเจน}
\index{การทำปุ๋ยหมัก (Composting Microbiology)}

\textbf{การทำปุ๋ยหมัก (Composting)} คือกระบวนการย่อยสลายทางชีวภาพของสารอินทรีย์โดยจุลินทรีย์หลากหลายชนิดภายใต้สภาวะที่มีออกซิเจนเพียงพอ เพื่อเปลี่ยนเศษซากพืชและมูลสัตว์ให้กลายเป็นสารปรับปรุงดินที่มีความเสถียรและอุดมไปด้วยฮิวมัส (Humus) ซึ่งปราศจากเชื้อโรคพืช เมล็ดวัชพืช และกลิ่นเหม็น

ในระหว่างกระบวนการหมักแบบใช้ออกซิเจน พลังงานเคมีที่ถูกปลดปล่อยออกมาจากการหายใจระดับเซลล์ของจุลินทรีย์จะเปลี่ยนเป็นพลังงานความร้อน ก่อให้เกิดการเปลี่ยนแปลงของอุณหภูมิและการสืบแทนของกลุ่มจุลินทรีย์อย่างเป็นลำดับ 4 ระยะ:
\begin{enumerate}
    \item \textbf{ระยะเมโซฟิลิกเริ่มต้น (Initial Mesophilic Phase: 20--45 $^\circ\text{C}$)} จุลินทรีย์กลุ่มชอบอุณหภูมิปานกลาง (แบคทีเรียและเชื้อรา) เจริญอย่างรวดเร็วโดยใช้น้ำตาลและกรดอะมิโนที่ย่อยง่าย ค่า pH อาจลดลงเล็กน้อยเนื่องจากการสร้างกรดอินทรีย์
    \item \textbf{ระยะเทอร์โมฟิลิก (Thermophilic Phase: 45--70 $^\circ\text{C}$)} เมื่ออุณหภูมิสูงเกิน 45 $^\circ\text{C}$ แบคทีเรียกลุ่มชอบความร้อนสูง เช่น \textit{Bacillus stearothermophilus}, \textit{Thermus} และแอกทิโนมัยซีททนร้อน จะเข้าแทนที่และเร่งย่อยสลายเซลลูโลสและเฮมิเซลลูโลส อุณหภูมิที่สูงกว่า 55 $^\circ\text{C}$ ติดต่อกันอย่างน้อย 3 วัน มีความสำคัญอย่างยิ่งในการทำลายเชื้อจุลินทรีย์ก่อโรคพืช ไข่พยาธิ และเมล็ดวัชพืชให้หมดไป
    \item \textbf{ระยะลดอุณหภูมิ (Cooling Phase: 45--35 $^\circ\text{C}$)} เมื่อสารอาหารที่ย่อยง่ายลดลง กิจกรรมของจุลินทรีย์ชะลอตัว อุณหภูมิจะลดลง จุลินทรีย์กลุ่มเมโซไฟล์และเชื้อราจะกลับเข้ามาเจริญอีกครั้งเพื่อย่อยสลายลิกนินที่เหลืออยู่
    \item \textbf{ระยะบ่มสุก (Maturation / Curing Phase)} อุณหภูมิลดลงเท่ากับอุณหภูมิบรรยากาศ แอกทิโนมัยซีทและเชื้อราสร้างสารประกอบฮิวมัสที่เสถียร ปุ๋ยหมักที่ได้จะมีสีน้ำตาลเข้ม กลิ่นคล้ายหน้าดินในป่า และพร้อมนำไปใช้
\end{enumerate}

\begin{figure}[H]
\centering
\begin{tikzpicture}[scale=1.0]
    % Axes
    \draw[-{Stealth[scale=1.2]}, line width=1.5pt, color=agriSlate] (0,0) -- (10.5,0) node[right, font=\footnotesize\bfseries] {เวลา (สัปดาห์)};
    \draw[-{Stealth[scale=1.2]}, line width=1.5pt, color=agriSlate] (0,0) -- (0,5.5) node[above, font=\footnotesize\bfseries] {อุณหภูมิ ($^\circ\text{C}$)};

    % Ticks
    \node at (0, 1.5) [left, font=\scriptsize] {30};
    \node at (0, 2.7) [left, font=\scriptsize] {45};
    \node at (0, 4.2) [left, font=\scriptsize] {65};
    \draw[dashed, color=slateGrey!40] (0, 2.7) -- (10, 2.7);
    \draw[dashed, color=red!40] (0, 4.2) -- (10, 4.2);

    % Pathogen kill threshold label
    \node at (5.0, 4.5) [font=\tiny\bfseries\color{agriRed}] {ขอบเขตทำลายเชื้อก่อโรคและเมล็ดวัชพืช ($>55^\circ\text{C}$)};

    % Composting curve
    \draw[line width=2.5pt, color=agriRed!80!black] 
        plot [smooth] coordinates {(0,1.2) (0.8, 2.8) (2.0, 4.8) (3.5, 4.5) (5.5, 2.6) (7.5, 1.8) (9.5, 1.4)};

    % Phase divisions
    \draw[dotted, line width=1pt, color=agriSlate] (1.2, 0) -- (1.2, 5.0);
    \draw[dotted, line width=1pt, color=agriSlate] (4.8, 0) -- (4.8, 5.0);
    \draw[dotted, line width=1pt, color=agriSlate] (7.2, 0) -- (7.2, 5.0);

    % Phase labels
    \node at (0.6, 5.2) [font=\tiny\bfseries\color{agriGreen}, align=center] {1. เมโซฟิลิก\\(Mesophilic)};
    \node at (3.0, 5.2) [font=\tiny\bfseries\color{agriRed}, align=center] {2. เทอร์โมฟิลิก\\(Thermophilic)};
    \node at (6.0, 5.2) [font=\tiny\bfseries\color{agriGold!80!black}, align=center] {3. ลดอุณหภูมิ\\(Cooling)};
    \node at (8.6, 5.2) [font=\tiny\bfseries\color{agriEarth}, align=center] {4. บ่มสุก\\(Curing)};
\end{tikzpicture}
\caption{กราฟแสดงการเปลี่ยนแปลงอุณหภูมิและระยะการย่อยสลายในกระบวนการทำปุ๋ยหมักแบบใช้ออกซิเจน}
\label{fig:compost_temperature}
\end{figure}

\section{การหมักแบบไม่ใช้ออกซิเจนและการผลิตก๊าซชีวภาพ}
\index{ก๊าซชีวภาพ (Biogas Production)}

\textbf{การย่อยสลายแบบไม่ใช้ออกซิเจน (Anaerobic Digestion: AD)} เป็นกระบวนการที่กลุ่มจุลินทรีย์ไม่ใช้ออกซิเจนทำงานประสานกันเป็นลูกโซ่ทางชีวเคมี 4 ขั้นตอน:
\begin{enumerate}
    \item \textbf{ไฮโดรไลซิส (Hydrolysis)} เอนไซม์ภายนอกเซลล์ย่อยสลายสารอินทรีย์โมเลกุลใหญ่ (คาร์โบไฮเดรต โปรตีน ลิพิด) ให้เป็นโมเลกุลเดี่ยว (กลูโคส กรดอะมิโน กรดไขมัน)
    \item \textbf{การสร้างกรด (Acidogenesis)} แบคทีเรียเปลี่ยนโมเลกุลเดี่ยวให้เป็นกรดไขมันระเหยง่าย (Volatile Fatty Acids: VFAs) เช่น กรดโพรพิโอนิกและกรดบิวทีริก
    \item \textbf{การสร้างกรดแอซีติก (Acetogenesis)} แบคทีเรียกลุ่มแอซีโตเจนเปลี่ยน VFAs ให้เป็นกรดแอซีติก ($CH_3COOH$), ก๊าซคาร์บอนไดออกไซด์ ($CO_2$) และไฮโดรเจน ($H_2$)
    \item \textbf{การสร้างก๊าซมีเทน (Methanogenesis)} อาร์เคียกลุ่มเมทาโนเจน (Methanogenic Archaea) เช่น \textit{Methanosarcina} และ \textit{Methanobacterium} เปลี่ยนกรดแอซีติกและไฮโดรเจนให้กลายเป็นก๊าซมีเทน ($CH_4$) ซึ่งเป็นส่วนประกอบหลักของก๊าซชีวภาพ (55--70\%)
\end{enumerate}

\begin{agribox}{แบคทีเรียสังเคราะห์แสงสีแดงกับการเพิ่มผลผลิตข้าว}
แบคทีเรียสังเคราะห์แสงกลุ่มสีม่วงและแดง (\textit{Rhodopseudomonas palustris}) สามารถใช้วัสดุอินทรีย์และก๊าซไฮโดรเจนซัลไฟด์ ($H_2S$) ในดินนาข้าวเป็นแหล่งอิเล็กตรอนในการสังเคราะห์แสง ช่วยลดกลิ่นเหม็นและลดความเป็นพิษของซัลไฟด์ต่อรากข้าว พร้อมทั้งสังเคราะห์กรดอะมิโน 5-aminolevulinic acid (ALA) ซึ่งเป็นสารกระตุ้นการสร้างคลอโรฟิลล์ ทำให้ต้นข้าวเขียวนาน แตกกอดี และเมล็ดเต่ง
\end{agribox}

\section*{คำถามทบทวนท้ายบทที่ 6}
\addcontentsline{toc}{section}{คำถามทบทวนท้ายบทที่ 6}
\begin{enumerate}
    \item เหตุใดการรักษาอุณหภูมิในระยะ Thermophilic ($>55^\circ\text{C}$) จึงเป็นข้อกำหนดมาตรฐานในการผลิตปุ๋ยหมักทางการค้า
    \item จงอธิบายบทบาทของจุลินทรีย์ใน 4 ขั้นตอนของการผลิตก๊าซชีวภาพ (Anaerobic Digestion)
    \item หากกองปุ๋ยหมักมีอัตราส่วน C:N ratio สูงเกินไป (เช่น มากกว่า 60:1) จะส่งผลกระทบต่ออัตราการย่อยสลายอย่างไร
\end{enumerate}

\clearpage